\chapter{绪论}

\section{研究背景与意义}

细粒度图像分类关注的不是“大类是否正确”,而是同一大类内部的子类区分。例如,普通分类任务只需判断图像是猫还是狗,细粒度分类则要求继续识别英短、布偶、暹罗、柴犬、哈士奇等具体品种。与粗粒度分类相比,该任务通常同时具有类间差异小和类内变化大的特点: 不同类别之间可能只在耳朵形状、毛发纹理、脸部轮廓或体态比例上存在局部差别,而同一类别又会因为姿态、光照、背景和遮挡的变化呈现出明显不同的外观\cite{parkhi2012cats,atabuzzaman2025zero}。这使得细粒度分类对模型的特征表达能力提出了更高要求。

从模型发展过程来看,深度学习显著推动了图像分类性能提升。AlexNet 在 ImageNet 上取得突破后,卷积神经网络开始成为视觉任务的基础方案\cite{krizhevsky2012imagenet}。随后,VGG 和 ResNet 等结构进一步改善了模型深度和训练稳定性\cite{simonyan2015very,he2016deep}。这些方法在标准分类任务上取得了很强的效果,但它们通常建立在大量标注数据的基础上。对于细粒度任务而言,高质量标注往往更难获取,一方面是因为类别数较多,另一方面是因为部分类别之间需要较强的领域知识才能准确区分。

宠物品种识别就是一个很典型的例子。Oxford-IIIT Pets 数据集包含 37 个猫狗品种,图像数量有限,并且样本在姿态、尺度、背景和拍摄质量上变化明显\cite{parkhi2012cats,oxfordpetsweb}。已有研究表明,即使是犬种分类这样接近日常应用的任务,卷积网络在面对外观相近类别时仍然容易产生混淆\cite{cuizhang2023classification}; 对于猫类识别,细粒度差异通常更依赖局部特征,迁移学习与分层特征利用才可能取得更稳定的结果\cite{wang2024enhancing}。因此,围绕宠物品种展开细粒度分类研究,不仅具有明确的实验价值,也具备实际应用意义。

近年来,视觉语言预训练模型为小样本视觉任务提供了新的技术路径。CLIP 通过大规模图文对比学习建立了统一的图像和文本语义空间,使模型可以借助文本提示完成零样本分类\cite{clip2021}。与传统分类网络相比,它不需要为每个任务单独训练完整分类头,在标注数据有限的场景下具有明显优势。不过,CLIP 在下游任务中的表现对提示词设计高度敏感。对于细粒度类别,如果提示词过于通用,文本语义往往不足以支撑不同品种之间的细微区分。

围绕这一问题,提示学习逐渐成为视觉语言模型适配的重要方向\cite{gu2023prompt,chen2023vlp}。研究者开始尝试用可学习上下文替代人工模板,并进一步引入图像条件提示和细粒度多模态提示,以改善 CLIP 在少样本场景中的表现\cite{coop2022,cocoop2022,liu2025fgmpl}。这说明提示词已经不再只是输入文本的修饰,而逐渐成为模型适配策略本身的一部分。

本文的工作正是在这一背景下展开。研究重点是: 在 Oxford-IIIT Pets 数据集的少样本设置下,如何通过动态提示优化提高 CLIP 对细粒度猫狗品种的分类能力。相比于直接微调整个模型,提示学习参数量更小,训练成本更低,也更适合本科毕业设计的实验条件; 相比于固定模板,动态提示则更有可能适应不同图像和不同分类难度的样本。

\section{国内外研究现状}

\subsection{细粒度图像分类研究}

细粒度分类并不是随着大模型才出现的研究方向。早期方法更多依赖人工设计特征、部件检测和属性建模,希望通过局部区域或相对属性来刻画子类别之间的差异。Oxford-IIIT Pets 数据集提出时,研究者就已经把宠物品种识别视作细粒度视觉识别问题,并指出该任务对局部形状和毛发表观建模提出了较高要求\cite{parkhi2012cats}。这类方法有助于解释模型关注了哪些区域,但往往依赖较强的任务先验和人工特征设计。

深度学习兴起后,细粒度分类逐渐转向端到端特征学习。卷积神经网络凭借层级表示能力,在鸟类、车辆和宠物等任务上取得了更稳定的结果。围绕宠物分类,已有研究尝试采用卷积网络结合传统分类器实现犬种识别,结果表明类别相似性和样本复杂度仍然是限制性能的重要因素\cite{cuizhang2023classification}。也有工作从迁移学习角度改进细粒度猫类识别,说明对中高层特征进行更细致的利用能够带来性能提升\cite{wang2024enhancing}。

随着视觉语言模型的发展,细粒度分类又出现了新的切入点。Atabuzzaman 等人在 EMNLP 2025 Findings 中讨论了大规模视觉语言模型在零样本细粒度分类上的潜力,指出更合理的类别描述和提示设计能够显著影响模型表现\cite{atabuzzaman2025zero}。这一结果说明,语言信息并不是附属输入,在细粒度任务中反而可能成为增强类别区分能力的重要补充。

\subsection{视觉语言预训练与提示学习研究}

CLIP 的提出使视觉语言预训练成为当前多模态研究的重要基础\cite{clip2021}。该模型通过大规模图文对学习跨模态对齐表示,图像特征和文本特征被映射到同一语义空间中,从而可以通过文本模板直接完成分类。围绕视觉语言预训练,后续研究和综述逐渐形成较为完整的知识框架。Chen 等人的综述系统梳理了视觉语言预训练的模型结构、训练目标和下游任务\cite{chen2023vlp}; Gu 等人则从提示工程角度总结了视觉语言基础模型中的提示形式与应用场景\cite{gu2023prompt}。这些工作都表明,提示已经成为视觉语言模型适配不可回避的核心问题。

在具体方法上,CoOp 是视觉任务提示学习的代表工作\cite{coop2022}。它将人工文本模板中的上下文词替换为可学习向量,并利用少量下游样本优化这些向量,从而使提示更符合目标数据集分布。CoOp 的关键贡献在于证明了一个事实: 在很多任务中,无需微调整个视觉编码器,只通过学习少量提示参数就可以显著提升 CLIP 的分类性能。

不过,CoOp 的提示是静态共享的,所有图像都使用同一组上下文向量。当任务内部样本变化较大时,这种统一提示很难同时适应所有输入。为此,Zhou 等人进一步提出 CoCoOp,在静态提示基础上增加图像条件 token,使不同图像能够获得不同的提示偏移\cite{cocoop2022}。CoCoOp 把提示学习从任务级共享推进到样本级适配,在少样本泛化和跨数据集测试中表现出更好的鲁棒性。

在 CoOp 和 CoCoOp 之后,提示学习开始向更复杂的方向发展。一类工作关注自动提示优化。ProAPO 通过渐进式策略对提示进行自动搜索和改进,以减少手工设计的参与程度\cite{qu2025proapo}。另一类工作强调多模态提示协同。Liu 等人在 Neurocomputing 2025 中提出细粒度多模态提示学习方法,同时考虑视觉和文本提示之间的交互,以增强视觉语言模型对细微类别差异的表达能力\cite{liu2025fgmpl}。此外,面向更广义多模态大模型的提示优化研究也表明,如果只关注单一模态提示,往往不能充分利用模型已有的跨模态对齐能力\cite{choi2025multimodal}。

\subsection{现有工作的不足}

综合现有研究可以看到,视觉语言模型和提示学习已经为细粒度分类提供了较强的技术基础,但仍有几个问题没有被充分解决。

第一,很多工作重点解决的是“如何生成更好的提示”,却较少直接讨论“训练时哪些样本更值得被强调”。细粒度任务中真正影响性能上限的往往不是那些很快就能分类正确的样本,而是处在类别边界附近、置信度较低的困难样本。如果训练时对所有样本一视同仁,模型很可能把较多更新分配给简单样本,而不是决定性能的难样本。

第二,已有条件提示方法虽然能够根据图像生成不同提示,但其提示调整主要依赖视觉特征本身,并没有把模型当前的预测状态显式纳入训练机制。换句话说,模型知道“图像内容是什么”,却未必知道“当前这张图究竟有多难分”。对细粒度分类而言,如果能把预测置信度和误分类情况引入损失加权,训练过程可能会更加贴合任务需求。

第三,自动提示优化和多模态提示方法虽然具有潜力,但不少方法结构较复杂、训练流程较长,对复现实验和工程实现有一定门槛\cite{qu2025proapo,liu2025fgmpl,choi2025multimodal}。对于以少样本宠物识别为目标的任务,仍需要一种结构清晰、参数量较小、同时能够体现样本难度差异的方法。

\section{本文主要研究内容}

针对以上问题,本文围绕基于提示词优化的细粒度猫狗分类方法展开研究,主要工作包括以下几个方面。

第一,分析 CLIP、CoOp 和 CoCoOp 的核心机制,梳理视觉语言预训练模型在细粒度少样本分类中的适用性与局限性,明确本文方法的切入点。

第二,设计并实现 DynamicPromptTrainer。该方法以 CLIP 为基础,通过 AdaptivePromptLearner、SoftPromptAdapter 和 DifficultyWeightCalculator 三个模块,实现类别自适应提示、图像条件提示和难度感知样本加权。

第三,在 Oxford-IIIT Pets 数据集上构建 1-shot、2-shot、4-shot、8-shot 和 16-shot 的实验设置,比较本文方法与 CoOp、CoCoOp 的分类表现,分析动态提示与样本加权对细粒度分类结果的影响。

第四,阅读并整合项目中的安全扩展模块,梳理基于 Test-time Counterattack 的推理时防御机制,说明其在系统中的集成方式与评估流程。

第五,结合项目实际训练过程,整理与 shot 数匹配的自适应 epoch 训练策略,并将其纳入最终实验方案,提高不同数据规模下实验对比的稳定性。

\section{本文组织结构}

全文共分为六章,结构安排如下。

第一章为绪论,介绍研究背景、国内外研究现状、本文研究内容和论文整体结构。

第二章为相关技术,重点介绍 CLIP、CoOp、CoCoOp 以及细粒度分类与难样本学习相关方法,为本文方法设计提供理论基础。

第三章为系统设计,详细说明 DynamicPromptTrainer 的整体架构与核心模块实现。

第四章为系统安全设计,介绍项目中加入的对抗性防御功能,说明其核心原理、代码结构、推理流程与评估方案。

第五章为实验与结果分析,给出数据集、实验设置、对比方法和最终结果,并结合图表分析本文方法在不同 few-shot 设置下的表现。

第六章为总结与展望,对全文工作进行总结,并讨论后续仍可继续改进的方向。
